% !TEX program = pdflatex
\documentclass[11pt,a4paper]{article}

\usepackage[utf8]{inputenc}
\usepackage[T1]{fontenc}
\usepackage[ngerman]{babel}
\usepackage{lmodern}
\usepackage{microtype}

\usepackage{amsmath,amssymb,amsthm,mathtools}
\usepackage{hyperref}

% --- "Serum-Sache": deutsch benannte Theorem-Umgebungen ---
\newtheorem{satz}{Satz}
\newtheorem{lemma}[satz]{Lemma}
\newtheorem{korollar}[satz]{Korollar}
\theoremstyle{definition}
\newtheorem{definition}[satz]{Definition}
\newtheorem{beispiel}[satz]{Beispiel}
\theoremstyle{remark}
\newtheorem{bemerkung}[satz]{Bemerkung}

\title{Spin-Erzwingung im Oktonen-Upgrade und Anschluss an Dualität $s\mapsto 1-s$}
\author{}
\date{}

\begin{document}
	\maketitle
	
	\section{Ziel und Idee in einem Satz}
	Wir formulieren eine \emph{präzise} ``Spin-/Zulässigkeitsregel'' für additive Verwendung
	von oktonischen Produkttermen und koppeln sie an eine Dualitäts-Involution $U$,
	die die funktionale Symmetrie $s\mapsto 1-s$ geometrisch repräsentiert.
	Das Ergebnis ist eine \emph{konditionale} Schlusskette:
	\[
	\text{(additive Stabilität)} \;\Rightarrow\; v(\rho)\in \mathrm{Fix}(U)
	\;\Rightarrow\; \Re(\rho)=\tfrac12,
	\]
	sofern das Embedding $v(s)$ die Fixmenge $\mathrm{Fix}(U)$ exakt mit der kritischen Achse
	identifiziert.
	
	\section{Oktonen: Assoziator und ``additive Stabilität''}
	\begin{definition}[Assoziator]
		In den reellen Oktonen $\mathbb O$ sei der \emph{Assoziator} definiert durch
		\[
		[x,y,z]\;:=\;(xy)z \;-\; x(yz).
		\]
	\end{definition}
	
	\begin{bemerkung}
		In $\mathbb H$ (Quaternionen) gilt $[x,y,z]\equiv 0$ (Assoziativität).
		In $\mathbb O$ ist $[x,y,z]$ im Allgemeinen nicht null (Nichtassoziativität).
		Die Oktonen sind jedoch \emph{alternativ}; insbesondere ist die von zwei Elementen erzeugte
		Unteralgebra assoziativ (``Artin-Eigenschaft''). Diese Beobachtung ist der Kern des
		``Quaternionen\,$\to$\,Oktonen''-Zwangs: additive Superposition ist stabil nur in
		assoziativen Sektoren.
	\end{bemerkung}
	
	\subsection{Interaktionskontext (die Menge der zulässigen ``dritten Faktoren'')}
	\begin{definition}[Interaktionsmenge]
		Sei $\mathcal Z\subset \mathbb O$ eine fest gewählte Menge von ``Interaktionsrichtungen''.
		In einem Toy-Modell kann man $\mathcal Z$ etwa als die Menge der relevanten Basis-Einheiten,
		oder als die Menge der in der Dynamik tatsächlich auftretenden Faktoren wählen.
	\end{definition}
	
	\begin{definition}[Additive Zulässigkeit (assoziatorbasiert)]
		Ein Produktterm $xy$ heißt \emph{additiv zulässig} (im Interaktionskontext $\mathcal Z$),
		falls er mit allen relevanten dritten Faktoren assoziiert:
		\[
		\boxed{\;\;xy\ \text{ist additiv zulässig}\;\Longleftrightarrow\;
			[x,y,z]=0\ \ \forall z\in\mathcal Z.\;\;}
		\]
	\end{definition}
	
	\begin{bemerkung}[Physikalische Lesart: ``Superpositionsfähigkeit'']
		Wenn $[x,y,z]\neq 0$ für ein relevantes $z$, dann hängt das Ergebnis späterer
		Dreifachprodukte (und damit die Dynamik) von der Klammerung ab.
		Eine lineare Superposition von Termen, deren spätere Wirkung klammerungsabhängig ist,
		ist als ``Zustandsaddition'' nicht robust.
		Die obige Definition ist eine präzise Form deiner Aussage:
		\emph{Produktterme dürfen nur dann additiv verwendet werden, wenn sie in einem
			quaternionischen (assoziativen) Sektor liegen.}
	\end{bemerkung}
	
	\section{Geometrische Kodierung: Winkelhalbierende als Fixraum einer Involution}
	Nun übersetzen wir ``oktonische Winkelhalbierende'' in eine Invarianzstruktur.
	
	\begin{definition}[Euklidischer Modellraum]
		Wir identifizieren den zugrunde liegenden reellen Vektorraum der Oktonen mit
		\[
		V:=\mathbb R^8
		\]
		mit dem Standard-Skalarprodukt $\langle \cdot,\cdot\rangle$ und Norm
		$\|v\|^2=\langle v,v\rangle$.
	\end{definition}
	
	\begin{definition}[Dualitätsinvolution und Fixraum (``Winkelhalbierende'')]
		Sei $U\in O(8)$ eine orthogonale Involution, also
		\[
		U^2=\mathrm{Id},\qquad \langle Uv,Uw\rangle=\langle v,w\rangle.
		\]
		Wir definieren die \emph{oktonische Winkelhalbierende} als Fixraum
		\[
		H\;:=\;\mathrm{Fix}(U)\;=\;\{v\in V:\ Uv=v\}.
		\]
	\end{definition}
	
	\begin{definition}[Additive Zulässigkeit (geometrisch)]
		Ein Produktterm $xy$ heißt \emph{geometrisch additiv zulässig}, falls sein Bild in $V$
		in $H$ liegt:
		\[
		\boxed{\;\;xy\ \text{zulässig}\;\Longleftrightarrow\; \iota(xy)\in H,\;\;}
		\]
		wobei $\iota:\mathbb O\to V$ die kanonische Identifikation ist.
	\end{definition}
	
	\begin{bemerkung}
		Die assoziatorbasierte und die geometrische Sicht lassen sich zusammenführen, indem man
		$U$ so wählt, dass $H$ genau den ``assoziativen/physikalischen'' Unterraum kodiert
		(z.\,B. als Fixraum einer fundamentalen Austausch-/Dualitätssymmetrie).
	\end{bemerkung}
	
	\subsection{Toy-Wahl: $b\leftrightarrow c$ als ``Spin-Flips'' und Fixpunktbedingung}
	Du hattest das Motiv: $ab\mapsto c$, $ac\mapsto b$, usw., aber additiv nur auf der
	Winkelhalbierenden.
	
	\begin{definition}[$b\leftrightarrow c$-Involution im Toy-Sektor]
		Wir wählen in einem dreidimensionalen Imaginärsektor $\mathrm{span}\{a,b,c\}\subset \mathbb O$
		eine Involution $U$ als linearen Austausch
		\[
		U(a)=a,\qquad U(b)=c,\qquad U(c)=b,
		\]
		und erweitern $U$ orthogonal auf ganz $V=\mathbb R^8$ (auf den übrigen Basisrichtungen
		als $\pm$-Identität, je nach Modell).
		Dann ist der Fixraum in diesem Sektor
		\[
		H\cap \mathrm{span}\{a,b,c\}=\mathrm{span}\Big\{a,\ \frac{b+c}{\sqrt2}\Big\}.
		\]
		Der Vektor $(b+c)/\sqrt2$ ist die ``Winkelhalbierende'' zwischen $b$ und $c$.
	\end{definition}
	
	\begin{bemerkung}[Dein Satz in Formeln]
		``$ab$ darf additiv nur auftreten, wenn es auf der Winkelhalbierenden liegt'' bedeutet:
		\[
		\iota(ab)\in \mathrm{span}\Big\{a,\frac{b+c}{\sqrt2}\Big\}
		\quad\text{bzw.}\quad
		U(\iota(ab))=\iota(ab).
		\]
		Analog für $ac$ und $bc$.
		Das ist eine konkrete, testbare Form des Spin-Zwangs.
	\end{bemerkung}
	
	\section{Anschluss an die Dualität $s\mapsto 1-s$ (Nullstellengeometrie)}
	Jetzt koppeln wir den oktonischen Spin-Zwang an die Zeta-/Xi-Dualität.
	
	\begin{definition}[Embedding der kritischen Streifenvariable]
		Sei
		\[
		v:\ \{s\in\mathbb C:0<\Re(s)<1\}\ \to\ V
		\]
		eine Abbildung (``Zustandsembedding'') in den oktonischen Modellraum.
	\end{definition}
	
	\begin{definition}[Dualitätskompatibilität]
		Wir nennen $v$ \emph{dualitätskompatibel} (bezüglich $U$), wenn für alle $s$ im kritischen
		Streifen gilt
		\[
		\boxed{\;\;v(1-s)=U\,v(s).\;\;}
		\]
	\end{definition}
	
	\begin{bemerkung}
		Das ist die geometrische Form der funktionalen Symmetrie $\xi(s)=\xi(1-s)$.
		Hier wird nicht $\xi$ selbst in $V$ abgebildet, sondern die Zustände/Punkte $s$.
		Diese Kopplung ist die zentrale Modellannahme.
	\end{bemerkung}
	
	\subsection{Spin-Zwang als Stabilitätsprinzip für Nullstellen}
	\begin{definition}[Nullstellen als additiv stabile Zustände (Spin-Prinzip)]
		Wir postulieren als \emph{Spin-/Stabilitätsprinzip}:
		Für jede nichttriviale Nullstelle $\rho$ der Xi-Funktion ist $v(\rho)$ ein additiv stabiler
		Zustand im oktonischen Sinn, d.\,h.
		\[
		\boxed{\;\;\xi(\rho)=0\ \Longrightarrow\ v(\rho)\in H=\mathrm{Fix}(U).\;\;}
		\]
	\end{definition}
	
	\begin{bemerkung}[Warum das genau dein ``Spin erzwingt'' ist]
		Das Prinzip sagt: Nullstellen sind keine beliebigen Punkte, sondern müssen in dem
		$U$-geraden (fixen) Sektor liegen, in dem additive Produktterme konsistent bleiben.
		Das ist eine mathematische Superselektionsregel:
		\emph{nur $U$-Eigenzustände mit Eigenwert $+1$ sind physikalisch/additiv zulässig.}
	\end{bemerkung}
	
	\section{Fixraum $\Rightarrow$ kritische Achse}
	Nun kommt der Schritt, der die Fixraum-Bedingung in die Aussage $\Re(\rho)=1/2$ übersetzt.
	Dazu braucht man eine Identifikationsannahme: dass $H$ exakt die kritische Achse kodiert.
	
	\begin{definition}[Achsen-Kompatibilität des Embeddings]
		Wir sagen, $v$ sei \emph{achsen-kompatibel}, wenn gilt:
		\[
		\boxed{\;\;v(s)\in H\ \Longleftrightarrow\ \Re(s)=\tfrac12.\;\;}
		\]
	\end{definition}
	
	\begin{lemma}[Fixraum erzwingt kritische Linie]
		Ist $v$ dualitätskompatibel und achsen-kompatibel, dann gilt für jedes $s$:
		\[
		v(s)\in H\ \Rightarrow\ \Re(s)=\tfrac12.
		\]
	\end{lemma}
	
	\begin{proof}
		Dies ist die direkte Richtung der Achsen-Kompatibilität.
	\end{proof}
	
	\begin{satz}[Konditionale RH aus Spin/Stabilität]
		Angenommen, es existieren
		\begin{itemize}
			\item ein dualitätskompatibles und achsen-kompatibles Embedding $v$,
			\item eine Dualitätsinvolution $U\in O(8)$ mit Fixraum $H=\mathrm{Fix}(U)$,
			\item und das Spin-/Stabilitätsprinzip $\xi(\rho)=0\Rightarrow v(\rho)\in H$.
		\end{itemize}
		Dann liegen alle nichttrivialen Nullstellen $\rho$ der Xi-Funktion auf der kritischen Linie:
		\[
		\Re(\rho)=\tfrac12.
		\]
	\end{satz}
	
	\begin{proof}
		Sei $\rho$ eine nichttriviale Nullstelle. Nach dem Spin-/Stabilitätsprinzip gilt
		$v(\rho)\in H$. Nach dem Lemma folgt daraus $\Re(\rho)=\tfrac12$.
	\end{proof}
	
	\section{Relation zum ``Hurwitz-Fixpunkt''}
	Dein früherer Hurwitz-Fixpunkt war (inhaltlich):
	\[
	\| \rho\|^2=\|1-\rho\|^2\ \Rightarrow\ \Re(\rho)=\tfrac12.
	\]
	Hier ist die Rolle von $U$:
	Wenn $U$ orthogonal ist und $v(1-s)=Uv(s)$, dann gilt automatisch
	\[
	\|v(1-s)\|=\|Uv(s)\|=\|v(s)\|.
	\]
	Die Norminvarianz ist somit \emph{Folge} der Dualitätsisometrie (nicht ihr Ersatz).
	Der neue Beitrag ist, dass die Fixraum-Bedingung $v(\rho)\in H$ aus einer
	\emph{Nichtassoziativitäts-/Superpositionslogik} motiviert wird:
	Additive Stabilität im Oktonen-Upgrade zwingt in den $U$-geraden Sektor.
	
	\section{Was als Nächstes zu zeigen wäre (klar abgegrenzt)}
	Für einen ``harten'' mathematischen Fortschritt braucht man nun eine konkrete Konstruktion
	oder Charakterisierung von $v(s)$, die (i) dualitätskompatibel und (ii) achsen-kompatibel ist,
	und die das Spin-/Stabilitätsprinzip nicht als bloßes Postulat stehen lässt, sondern
	aus einem Positivitäts- oder Extremalitätszertifikat herleitet.
	
	Eine natürliche Zielrichtung (kompatibel mit E8/R8-Ideen) ist:
	\begin{itemize}
		\item $v(s)$ so zu wählen, dass $U$ die geometrische Umsetzung der Symmetrie
		$s\mapsto 1-s$ ist,
		\item und ``Stabilität'' als Extremalität eines dualitätsinvarianten Funktionals
		zu formulieren, das außerhalb von $H$ Nichtassoziativitätskosten (Assoziator-Energie)
		misst.
	\end{itemize}
	
\end{document}